\scanpage{23}
% archive sheet number: 22
\mspage{-II- 13 -}
l'adhérence d'une face de $g'(C)$. De plus, si $h(W_i)=g'(W_k)$, on a $r_p h(W_i)=r_p g'(W_k)$, d'où, par $(\mathrm{II}_{m-1})$, $(r_p h)r_i(r_p h)^{-1}=(r_p g')r_k(r_p g')^{-1}$ ; ceci donne aussitôt $hr_i h^{-1}=g'r_k g'^{-1}$. On a donc démontré $(\mathrm{II}_m)$ dans le second cas.

LEMME 5.- \uline{Soit} $g=r_{i_1}\ldots r_{i_m}\in G$, \uline{et soit} $W_j$ \uline{un mur}. \uline{Si les chambres} $C$ \uline{et} $g(C)$ \uline{sont de part et d'autre de} $W_j$, \uline{alors} $r_j g$ \uline{est un produit de} $m-1$ \uline{générateurs} $r_{i_k}$.

Considérons, en effet, la suite de chambres $C_k=r_{i_1}\ldots r_{i_k}(C)$ $(k=0,\ldots,m)$ ; on a $C_{\unsure{o|0}}=C$ et $C_m=g(C)$. Il existe deux éléments consécutifs de cette suite qui sont de part et d'autre de $W_j$, soient $C_{q-1}$ et $C_q$. Posons $h=r_{i_1}\ldots r_{i_{q-1}}$, de sorte que $C_{q-1}=h(C)$ et $C_q=hr_{i_q}(C)$. Comme $C$ et $r_{i_q}(C)$ ont en commun la face $W_{i_q}\cap\overline C$ de dimension $n-1$, $C_{q-1}$ et $C_q$ ont en commun une face $F$ de dimension $n-1$ contenue dans $h(W_{i_q})$. On a $F\subset\overline C_q\cap\overline C_{q-1}$, d'où $F\subset W_j$ puisque $C_{q-1}$ et $C_q$ sont de part et d'autre de $W_j$ ; comme $F$ engendre l'hyperplan $h(W_{i_q})$, on a donc $W_j\supset h(W_{i_q})$ et $W_j=h(W_{i_q})$ (puisqu'il s'agit de deux \unsure{hyperplans|hyperplens}). Le lemme 4 montre qu'on a $r_j=hr_{i_q}h^{-1}$, d'où
\[
r_jg=r_jhr_{i_q}r_{i_{q+1}}\ldots r_{i_m}
=hr_{i_q}r_{i_q}r_{i_{q+1}}\ldots r_{i_m}
=r_{i_1}\ldots r_{i_{q-1}}r_{i_{q+1}}\ldots r_{i_m},
\]
ce qui démontre notre assertion.

LEMME 6.- \uline{Soit} $g\in G$. \uline{Si} $g(C)\cap C\ne\varnothing$, \uline{alors} $g=1$.

En vertu du lemme 3, $g(C)$ ne rencontre aucun $W_i$ ; donc, si $g(C)\cap C\ne\varnothing$, on a $g(C)\subset C$ puisque $g(C)$ est connexe ; de même on a $g^{-1}(C)\subset C$, d'où $C\subset g(C)$ et $C=g(C)$. Supposons $g\ne1$, et exprimons-le comme un produit $g=r_{i_1}\ldots r_{i_m}$ de longueur minimum. Les chambres $C$ et $r_{i_1}(C)=r_{i_2}\ldots r_{i_m}(C)$ sont situées de part et d'autre de $W_{i_1}$ ; donc, d'après le lemme 5, $r_{i_1}(r_{i_2}\ldots r_{i_m})=g$ peut s'écrire sous la forme d'un
